\scanpage{77}
Pour $i$ grand $G,G'$ proviennent de $G_i,G'_i$ \inserted{sch. en grpes comm. de t.f. plats sur $X_i$}.
\deleted{et $\xi$ d'une corr.}
$G_i|U_i=A_i$ et $G'_i|U_i=A'_i$ des sch. abéliens, $\xi$ provient d'un élém de corr. div. $\xi'$ sur $A_i\times A'_i$ ; $T$ provient d'un tore $T_i\subset G'_{iY_i}$, $S$ provient d'un sous-schéma plat $S_i\subset G_{iY_i}$.
De plus, \deleted{dans le cas} \deleted{\unsure{Encore}} on peut \deleted{\unsure{sans restriction}} supposer $S$ connexe \inserted{lisse} (sinon le remplacer par $S^0$) et \deleted{\unsure{de la partie unipotente} \unsure{Comme si} $\operatorname{car}k>0$ \unsure{comme le quotient par} $\unsure{p^nS}$} extension \deleted{d'une} successive d'une V.A., d'un tore \deleted{\unsure{pour }$S$\unsure{ propre et de un grpe}} grpe un. On peut donc supposer le $=$ dévissage valable pour $S_i$.
\deleted{\unsure{Cela implique}}
On peut supposer enfin (remplaçant au besoin les $A_i$ par $A_i[1/\ell]$) que $\ell$ inv. sur les $X_i$, donc les $Y_i$. Mais les ${}_{\ell^\nu}S_i$ sont encore \unsure{finis|fini} sur $Y_i$. Bref, ttes les données proviennent de données analogues, avec un \deleted{$\unsure{X}$} $X_0$ \inserted{$=\operatorname{Spec}A_0$} \uline{de type fini} sur $\operatorname{Spec}\mathbb Z$.

De plus, si $\mathfrak p=\mathfrak m\cap A_0$, alors $A_0\mathbin{-}V$ induit $\widehat A_0\longrightarrow V$, \deleted{\unsure{comme}} Donc, on peut remplacer les mots ``$A_0$ de type fini sur $\operatorname{Spec}\mathbb Z$'' par ``$A_0$ $\mathfrak p$-adique séparé et complet, avec $A_0/\mathfrak p$ de t.f. sur $\operatorname{Spec}\mathbb Z$, et $A_0\mathbin{-}V$ continu pour les topologies adiques''. On est ramené alors au th. pour cet $X_0,Y_0,G_0,G'_0,\xi_0,S_0,T_0$, et il relève du cas 1.
\hfill CQFD.
\mspage{76}
